% \documentclass[12pt]{article}
\documentclass[
  reprint,
  12pt,
  onecolumn,
  nofootinbib,
]{revtex4-2}
\usepackage{fullpage}

\usepackage{amsmath,amssymb,amsfonts}

\usepackage[english]{babel}
\usepackage[bbgreekl]{mathbbol}
\usepackage{amsfonts}
\usepackage{bm}
\usepackage{mathtools}
\usepackage{graphicx}

\usepackage[
  usenames,
  dvipsnames,
  svgnames,
  table
]{xcolor}
\usepackage[
  colorlinks=true,
  citecolor=NavyBlue,
  linkcolor=NavyBlue,
  urlcolor=PaleGreen4
]{hyperref}
\usepackage[capitalise]{cleveref}

\DeclareSymbolFontAlphabet{\mathbb}{AMSb}
\DeclareSymbolFontAlphabet{\mathbbl}{bbold}

%% \newcommand{\Tmat}[1]{\ensuremath{\mathbbl{#1}}}
%% \newcommand{\Tvec}[1]{\overset{\raisebox{-2pt}{\text{\tiny$\bm\rightarrow$}}}{\mathbbl{#1}}}
\newcommand{\Marrow}[1]{\overset{\text{\tiny$\bm\leftrightarrow$}}{#1}}
\newcommand{\mat}[1]{\ensuremath{\mathbf{#1}}}
\newcommand{\Tmat}[1]{\ensuremath{\mathbbl{#1}}}
\newcommand{\Tvec}[1]{\overset{\raisebox{-2pt}{\text{\tiny$\bm\rightarrow$}}}{\mathbbl{#1}}}
\newcommand{\Marrowbbl}[1]{\overset{\raisebox{-1pt}{\text{\tiny$\bm\Leftrightarrow$}}}{#1}}
\newcommand{\Dmat}[1]{\ensuremath{\Marrowbbl{\mathbf{#1}}}}
\newcommand{\Mrarrowbbl}[1]{\overset{\raisebox{-1pt}{\text{\tiny$\bm\Rightarrow$}}}{#1}}
\newcommand{\Dvec}[1]{\ensuremath{\Mrarrowbbl{\mathbf{#1}}}}


\newcommand{\rme}{\mathrm{e}}
\newcommand{\rmi}{\mathrm{i}}
\newcommand{\rmd}{\mathrm{d}}

\newcommand{\ti}{\text{i}}
\newcommand{\ta}{\text{a}}
\newcommand{\te}{\text{e}}
\newcommand{\ts}{\text{s}}
\newcommand{\tF}{\text{f}}
\newcommand{\tH}{\text{h}}

\newcommand{\pha}{\Omega L_\ta/c}
\newcommand{\phs}{\Omega L_\ts/c}
\newcommand{\rsp}{\mathfrak{r}_{\ts+}}
\newcommand{\rsm}{\mathfrak{r}_{\ts-}}
\newcommand{\rso}{\mathfrak{r}_{\ts0+}}
\newcommand{\drp}{\mathfrak{\delta r}_{\ts+}}
\newcommand{\drm}{\mathfrak{\delta r}_{\ts-}}

\DeclareMathOperator{\re}{Re}
\DeclareMathOperator{\im}{Im}

\graphicspath{{../models/tresults}}

\begin{document}

Using the following sign conventions
\begin{itemize}
\item Fourier transforms are
  $a(\omega) = \int \rmd t\, \rme^{-\rmi\omega t} a(t)$
  so that propagation in time is a phase delay $\rme^{-\rmi\omega t}$.
\item Increasing a propagation length by a distance $x$ corresponds to a negative phase rotation $\Tmat{R}(-kx)$.
\item Reflection from the front (HR) surface of an optic is $-r$, reflection from the back (AR) surface of an optic is $+r$, and transmission is $+t$.
\item Positive motion of an optic perpendicular to its face corresponds to the front (HR) surface moving forward away from the optic, i.e.\ pushing on the back of a mirror produces a positive motion and radiation pressure incident on the front of a mirror produces a force in the negative direction.
\end{itemize}

\section{Phase modulation by moving mirror}

Let $\Tvec{A}_{\text{fi}}, \Tvec{A}_{\text{fo}}, \Tvec{A}_{\text{bi}}$, and $\Tvec{A}_{\text{bo}}$ be the fields incident on the front, reflected from the front, incident on the back, and reflected from the back respectively. Further let $\Tvec{A}_{0\alpha}$ be the DC part of the field $\alpha$ and $\Tvec{a}_\alpha$ be the AC part so that $\Tvec{A} = \Tvec{A}_{0\alpha} + \Tvec{a}_\alpha$. Now suppose the optic is moved a distance $x$ in the positive direction. This corresponds to a roundtrip shortening of the distance between a source and the front of the mirror by $2x$ and a roundtrip increase of the distance between a source and the back surface of the mirror by $2x$. The fields are thus
\begin{subequations}
\begin{align}
  \Tvec{A}_\text{fo} &= -r \Tmat{R}(+2kx) \Tvec{A}_\text{fi} + t\Tvec{A}_\text{bi} \\
  \Tvec{A}_\text{bo} &= t\Tvec{A}_\text{fi} + r\Tmat{R}(-2kx) \Tvec{A}_\text{bi}.
\end{align}
\end{subequations}
Noting that $\Tmat{R}(\theta) = \Tmat{1} + \Tmat{R}(\pi/2)\theta + O(\theta^2)$, the AC part is
\begin{subequations}
\begin{align}
  \Tvec{a}_\text{fo} &= - \frac{4\pi r}{\lambda} \Tmat{R}(\pi/2) \Tvec{A}_\text{0fi}\, x \\
  \Tvec{a}_\text{bo} &= - \frac{4\pi r}{\lambda} \Tmat{R}(\pi/2) \Tvec{A}_\text{0bi}\, x.
\end{align}
\end{subequations}
Note that the signs are the same.

\section{Radiation pressure on the front of a mirror}

The radiation pressure force is $F_\text{rp} = (P_\text{bi} + P_\text{bo} - P_\text{fi} - P_\text{fo}) / c$ where $P_\alpha = \Tvec{A}_\alpha^\dag \Tvec{A}_\alpha$ is the power of field $\alpha$. First consider for simplicity a laser incident only on the front of a mirror as shown in \cref{fig:HRMirrorRP_graph}. Since the AC contribution is $2\Tvec{A}_{0\alpha}^\dag \Tvec{a}_\alpha$, the AC parts are
\begin{subequations}
  \label{eq:HRMirrorRP}
\begin{align}
  fr.Fq.i &\equiv \Tvec{f}_\text{i}^\dag = - \frac{2}{c} \Tvec{A}_\text{0fi}^\dag \\
  fr.Fq.o &\equiv \Tvec{f}_\text{o}^\dag = - \frac{2}{c} \Tvec{A}_\text{0fo}^\dag \\
  fr.px &\equiv \Tvec{x} = - \frac{4\pi r}{\lambda} \Tmat{R}(\pi/2) \Tvec{A}_\text{0fi} \\
  EX.chi &= \chi
\end{align}
\end{subequations}
The conversion of force to mirror motion is just multiplication by the mechanical susceptibility $\chi$. An attempt to implement this made in \texttt{test\_sflu\_HRMirrorRP} in test\_optical\_spring.py.

\begin{figure}
  \centering
  \includegraphics[scale=0.9]{plot_HRMirrorRP_graph/HRMirrorRP_graph.pdf}
  \caption{Graph of a Fabry-Perot cavity using an ETM with only an HR surface. The graph edges are given in \cref{eq:HRMirrorRP}. Produced by \texttt{plot\_HRMirrorRPReduced\_graph} in test\_optical\_spring.py.}
  \label{fig:HRMirrorRP_graph}
\end{figure}


Radiation pressure forms a loop and the Gaussian elimination can done by hand as shown in \cref{fig:HRMirrorRPReduced} with the result
\begin{subequations}
\label{eq:HRMirrorRPReduced}
\begin{align}
  \Tmat{F} &= \left(\Tmat{1} - \Tvec{x}\chi\Tvec{f}_\text{o}^\dag\right)^{-1} \\
  fr.r &= \Tmat{F}\left(-r\Tmat{1} + \Tvec{x}\chi\Tvec{f}_\text{i}^\dag\right) \\
  fr.xq &=
  \chi \left[\Tvec{f}_\text{i}^\dag + \Tvec{f}_\text{o}^\dag\Tmat{F}
    \left(-r\Tmat{1} + \Tvec{x}\chi\Tvec{f}_\text{i}^\dag\right)\right] \\
  fr.xF &= \left(\Tmat{1} + \chi \Tmat{f}_\text{o}^\dag \Tmat{F} \Tvec{x}\right) \chi \\
  fr.px &= \Tmat{F}\Tvec{x} \\
  fr.pF &= \Tmat{F}\Tvec{x}\chi \\
  fr.xx &= \Tmat{1} + \chi \Tvec{f}_\text{o}^\dag \Tmat{F}\Tvec{x}.
\end{align}
\end{subequations}
This is successfully implemented for an optical spring demonstrating both the optical and mechanical effects of radiation pressure; a comparison with Optickle is shown in \cref{fig:optickle_compare}.

\begin{figure}
  \centering
  \includegraphics[scale=0.9]{plot_HRMirrorRPReduced_graph/HRMirrorRPReduced_graph.pdf}
  \caption{Graph of a Fabry-Perot cavity using an ETM with only an HR surface, as in \cref{fig:HRMirrorRP_graph}, but where the Gaussian elimination of the radiation pressure loop has already been done. The graph edges are given in \cref{eq:HRMirrorRPReduced}. Produced by \texttt{plot\_HRMirrorRPReduced\_graph} in test\_optical\_spring.py.}
  \label{fig:HRMirrorRPReduced}
\end{figure}

\begin{figure}
  \centering
  \includegraphics[width=0.45\paperwidth]{test_HRMirrorRPReduced_compare_optickle/optical_tf.pdf}
  \includegraphics[width=0.45\paperwidth]{test_HRMirrorRPReduced_compare_optickle/mechanical_tf.pdf}
  \caption{Comparison of SFLU model of an optical spring given in \cref{eq:HRMirrorRPReduced} and \cref{fig:HRMirrorRPReduced} with an Optickle model of the same setup. Produced by \texttt{test\_HRMirrorRPReduced\_compare\_optickle} in test\_optical\_spring.py. (Same plots not requiring Optickle are made by \texttt{test\_sflu\_HRMirrorRPReduced}).}
  \label{fig:optickle_compare}
\end{figure}

\section{Radiation pressure on both sides of a mirror}

This is similar to above except there are now two radiation pressure loops which are more annoying to reduce by hand.

\end{document}
